\section{Camino m\'inimo rodeando un pol\'igono convexo}
\whenuse{El segmento directo est\'a bloqueado por un \`unico obst\'aculo convexo y
se permite recorrer su borde entre puntos tangentes.}

Usar cuando el segmento directo entre $A$ y $B$ atraviesa el interior de un
obst\'aculo convexo. Si no lo atraviesa, la respuesta es $|AB|$. Si lo atraviesa,
un camino \`optimo toca dos v\'ertices tangentes y recorre una de las dos cadenas
del borde entre ellos.

Procedimiento:
\begin{enumerate}
  \item hallar desde $A$ y $B$ los dos v\'ertices tangentes al pol\'igono;
  \item precomputar longitudes prefijas del per\'imetro;
  \item probar las combinaciones compatibles de tangencias y ambas cadenas;
  \item minimizar $|A-u|+\operatorname{borde}(u,v)+|v-B|$.
\end{enumerate}

\complexity{tangencias $O(\log n)$ con b\'usqueda binaria, o $O(n)$ en una
implementaci\'on simple; memoria $O(n)$ para prefijos.}

\warning{Si se permite tocar el borde, distinguir intersecci\'on con el interior
de un simple contacto tangencial. Para varios obst\'aculos construir un grafo de
visibilidad; esta receta aislada ya no alcanza.}
